Formalization is in progress. The shared library defines actual determinant-one matrices $n(x)$, $a(y)$, and $k(\theta)$ with the manuscript's entries. The companion proves their elementary actions, upper-half-plane transitivity, the full stabilizer as rotation matrices, existence of an $NAK$ factorization, and uniqueness of $x$, $y$, and the rotation matrix. It also proves the entry equations for the stabilizer, the inverse action of $s(x+iy)$, the explicit $x,y$ coordinates of $g\cdot i$, and the orbit-map criterion $gi=hi$ if and only if $h^{-1}g$ lies in the stabilizer.

The separated manuscript results are the \hyperref[hyperbolic.stabilizer]{stabilizer description}, \hyperref[hyperbolic.iwasawa]{Iwasawa factorization}, \hyperref[hyperbolic.iwasawa-explicit-coordinates]{explicit coordinates}, and \hyperref[hyperbolic.quotient-description]{quotient description}. The companion headings identify which of these results each existing declaration supports; splitting the prose does not extend its formal coverage.

The factorization proof follows the prose: write $g\cdot i=x+iy$; the position matrix $s(x+iy)$ sends $i$ to this same point, so its inverse times $g$ fixes $i$. The inverse-action theorem expresses the two undoing steps as $(w-x)/y$. The formal stabilizer proof compares real and imaginary parts, giving $a=d$ and $b=-c$, and then chooses an angle using the complex argument of $a+bi$. The coordinate proof uses Mathlib's complex division formulas, which encode multiplication by the conjugate denominator. The uniqueness proof applies both factorizations to $i$, compares real and imaginary parts, and cancels the common invertible factor. The square-root and rotation-sign conventions agree with the manuscript.

Equality of the rotation matrices is proved, but its reformulation as equality of real angles modulo $2\pi$ is not yet formalized. The explicit sine and cosine coordinates, identification with Mathlib's orthogonal-group interface, the quotient homeomorphism and smooth coordinates, and the angular derivative and double-covering proposition remain open. The Section 1.3 shared library proves the derivative at the single angle $\pi/4$ needed for metric uniqueness; the full angular formula and covering statement still require proof. The algebraic orbit criterion alone does not prove the quotient homeomorphism. Haar measure is treated in Section 1.8 and is outside this section's formalization.

Machine checks of the implemented portions are recorded in \texttt{lean/audits/chapter01-release-check-2026-11-12/record.json}, against its source hashes. Steven's correspondence review remains pending; the section remains in progress.
